\paragraph{Finite-field shadows and Jordan hairs: eventual polynomial law of recurrence dynamics in resonance windows and computable period certificates}
Once the minimal integer-coefficient recurrence in the resonance window has been audited and fixed (Theorem~\ref{thm:pom-moment-resonance}), the reductions and factor exponents of the characteristic polynomial \(P_q(\lambda)\) over finite fields \(\mathbb{F}_p\) compress the long-term dynamics of the sequence \(S_q(m)\) into a suite of completely internal, replicable fingerprint interfaces: they depend neither on additional continuous-limit hypotheses nor on external approximate models but are given directly by computable finite-dimensional algebraic objects. In connection with the modal-resonance chains in the main text (Remarks~\ref{rem:pom-charpoly-mod2-shadow},~\ref{rem:pom-charpoly-mod3-selection}), this class of interfaces can be summarized in the following five complementary assertions.
\begin{enumerate}
  \item \textbf{Eventual polynomial law modulo \texorpdfstring{$2$}{2}: difference annihilation compresses parity dynamics to finite-order discrete polynomials.}
  In the resonance window \(q=9,\dots,17\), the modulo-\(2\) shadow collapses to
  \(P_q(\lambda)\equiv \lambda^{a_q}(\lambda+1)^{b_q}\ (\mathrm{mod}\ 2)\), so the shift operator satisfies
  \(E^{a_q}(E+1)^{b_q}S_q\equiv 0\ (\mathrm{mod}\ 2)\).
  Writing \(\mathsf{D}:=E+1\) (identical to forward difference in characteristic \(2\)), we have in the recurrence-effective regime
  \(\mathsf{D}^{b_q}S_q(m+a_q)\equiv 0\ (\mathrm{mod}\ 2)\), so \(S_q(m+a_q)\bmod 2\) degenerates in the long term to a discrete polynomial of degree \(<b_q\) (Remark~\ref{rem:pom-charpoly-mod2-shadow}; Table~\ref{tab:fold_collision_charpoly_mod2_shadow_q9_17}).
  Thus \((a_q,b_q)\) are not only formal factorization exponents but also serve as complete invariants for ``eventual polynomial order/transient order,'' characterizing the deterministic degenerate structure of parity dynamics in the sufficiently late segment.
  In the lower-order verified window \(4\le q\le 11\), the exponent pair further condenses to the closed form
  \(b_q=2\lfloor (q+1)/4\rfloor\), \(a_q=3+2\,\ind\{q\ \mathrm{even},\ q\ge 6\}\), exhibiting a \(q\mapsto q+4\) staircase phase (Proposition~\ref{prop:pom-charpoly-mod2-closed-exponent-law-q4-11}).

  \item \textbf{Two-sector decomposition modulo \texorpdfstring{$3$}{3}: direct-sum compression of constant-phase and alternating-phase polynomials.}
  When \(P_q(\lambda)\equiv \lambda^2(\lambda-1)^4(\lambda+1)^5\ (\mathrm{mod}\ 3)\) (e.g., \(q=13,15\); Remark~\ref{rem:pom-charpoly-mod3-selection}), for \(\widetilde S_q:=E^2S_q\) we have
  \((E-1)^4(E+1)^5\widetilde S_q\equiv 0\ (\mathrm{mod}\ 3)\), so there exist two classes of generalized-phase modes controlled by the exponent pair \((4,5)\): one class consists of polynomials in \(m\) of degree \(\le 3\), the other class is \((-1)^m\) times polynomials of degree \(\le 4\) (Proposition~\ref{prop:pom-moment-mod3-two-phase}; Table~\ref{tab:fold_collision_charpoly_mod3_selection_q13_15}).
  This two-sector decomposition compresses all complexity in the modulo-\(3\) world into two blocks ``constant phase + alternating phase,'' providing foundational constraints for subsequent Jordan defect detection and embedding event localization.

  \item \textbf{\texorpdfstring{$A_2$}{A2} embedding and discriminant vanishing juxtaposition: local spectral bifurcation points as finite-field diagnostics.}
  When \(\chi_{A_2}(\lambda)\mid P_q(\lambda)\) holds in \(\mathbb{F}_p[\lambda]\), a three-dimensional invariant sub-dynamics appears in the modulo-\(p\) recurrence state space whose companion action has characteristic polynomial precisely \(\chi_{A_2}\) (Proposition~\ref{prop:pom-charpoly-modp-a2-embedding}; Corollary~\ref{cor:pom-charpoly-modp-a2-subdynamics}).
  In the verified slices \((q,p)=(14,7)\), \((12,11)\), this embedding flag co-occurs with \(\mathrm{repeat?}=1\), so \(\widetilde P_q:=P_q/\lambda^{a_0}\) is non-square-free and satisfies \(\gcd(\widetilde P_q,\partial_\lambda\widetilde P_q)\neq 1\); equivalently, its discriminant vanishes modulo \(p\) (Remark~\ref{rem:pom-a2-embedding-local-ramification}).
  Thus embedding events can be identified as local spectral bifurcation points accompanied by Jordan modulation channels, forming a joint fingerprint ``subrepresentation embedding \(\times\) ramification/non-semisimplicity.''

  \item \textbf{Discriminant blacklist \texorpdfstring{$\mathcal{B}_q$}{Bq} and global permutation symmetry: complete operationalization of exceptional-prime sets and testable frequency laws.}
  At \(q=9,10,11,13\), the observable minimal polynomial \(P_q\) first displays fixed-genus degree collapse while remaining irreducible (Proposition~\ref{prop:pom-observable-minpoly-order-collapse-q9-13}), and its integer discriminant has auditable structured factorizations (equation~\eqref{eq:fold_collision_observable_minpoly_discriminants_q9_13}).
  Thus the ramification prime blacklist
  \(\mathcal{B}_q=\{p:\ p\mid \mathrm{Disc}(P_q)\}\)
  can be directly generated from \(\mathrm{Disc}(P_q)\), and for any \(p\notin\mathcal{B}_q\) local multiple roots/spectral ramification are strictly excluded (Proposition~\ref{prop:pom-observable-minpoly-discriminant-ramification-q9-13}), elevating ``exceptional prime set'' to a completely finite auditable set.
  Furthermore, on the same slice \(\mathrm{Gal}(P_q/\QQ)\) attains the maximal symmetric group \(S_d\) (Proposition~\ref{prop:pom-observable-minpoly-galois-sd-q9-13}), so Chebotarev density yields a strict frequency law for splitting types, in particular the irreducibility density is \(1/d\) (Corollary~\ref{cor:pom-observable-minpoly-chebotarev-density-q9-13}); and Frobenius parity can be read out directly from the Kronecker symbol \(\left(\frac{\Delta_q^{\mathrm{disc}}}{p}\right)\) of the discriminant square class (Corollary~\ref{cor:pom-observable-minpoly-frob-parity-q9-13}).

  \item \textbf{Unified formula for period bound and criterion for extra \texorpdfstring{$p$}{p} factor: Jordan non-semisimplicity detectable without full factorization.}
  Suppose
  \(P_q(\lambda)\equiv \lambda^{a_0}\prod_j f_j(\lambda)^{e_j}\ (\mathrm{mod}\ p)\), where \(f_j\neq \lambda\) irreducible, \(d_j=\deg f_j\), and let \(\varepsilon_{q,p}=\ind\{\exists j:\ e_j\ge 2\}\).
  Then the coarse period upper bound in this paper's tables is uniformly given by
  \(
  \mathrm{PerBound}_p(q)=p^{\varepsilon_{q,p}}\cdot \mathrm{lcm}_j(p^{d_j}-1)
  \)
  (Remark~\ref{rem:pom-charpoly-modp-audit-method}; Tables~\ref{tab:fold_collision_charpoly_mod7_selection_q13_15}--\ref{tab:fold_collision_charpoly_mod13_fingerprint_q13}).
  The extra factor \(p\) equivalently characterizes non-square-freeness via \(\mathrm{repeat?}=1\), thus equivalently characterizing Jordan/generalized-eigenvector contributions in finite-field recurrence; and \(\varepsilon_{q,p}\) can be read out directly via \(\gcd(\widetilde P_q,\partial_\lambda\widetilde P_q)\) without requiring full factorization of \(P_q\).
\end{enumerate}

